\documentclass[12pt]{article}
\usepackage{amsmath}
\usepackage{amssymb}
\usepackage{geometry}
\geometry{a4paper, margin=1in}

\begin{document}

\section*{Exercise 1: Dynamic Table (INSERT and DELETE)}

We analyze a dynamic table $T$ that supports both \texttt{INSERT} and \texttt{DELETE} operations. To maintain $O(1)$ amortized cost, we use the standard strategy:
\begin{itemize}
    \item \textbf{Expansion:} Double the table size when inserting into a full table (load factor $\alpha = 1$).
    \item \textbf{Contraction:} Halve the table size when the table becomes less than 1/4 full ($\alpha = 1/4$).
\end{itemize}
We will prove the amortized cost per operation is $O(1)$ using all three methods.

\subsection*{1. Aggregate Analysis}
In aggregate analysis, we compute the total cost of a sequence of $n$ operations and divide by $n$.
\begin{itemize}
    \item The basic \texttt{INSERT} or \texttt{DELETE} without resizing costs $1$. For $n$ operations, this contributes at most $n$ to the total cost.
    \item \textbf{Cost of Expansions:} The table doubles at sizes $1, 2, 4, 8, \dots, 2^k$. The total cost of copying elements during expansions over $n$ operations is bounded by the geometric series: $\sum_{i=0}^{\lfloor \log n \rfloor} 2^i \le 2n$.
    \item \textbf{Cost of Contractions:} The table halves only when $\alpha$ drops below $1/4$. The cost of copying elements during contractions is similarly bounded by $O(n)$ because we must perform at least $\Theta(size(T))$ cheap deletes to trigger a contraction.
\end{itemize}
uTotal actual cost for $n$ operations: $\le n + 2n + n = 3n= O(n)$. \\
The amortized cost per operation is $\frac{O(n)}{n} = \mathbf{O(1)}$.

\subsection*{2. Accounting Method}

Assign the following amortized charges:

\[
\hat{c}_{\text{INSERT}} = 3
\qquad\text{and}\qquad
\hat{c}_{\text{DELETE}} = 2
\]

We show that these charges are sufficient to pay for every actual cost.

\subsubsection*{INSERT}

For each INSERT:

\begin{itemize}
    \item 1 unit pays for the immediate insertion,
    \item 2 units are stored as credit with the inserted element.
\end{itemize}

Suppose an expansion occurs when the table has capacity $k$ and currently contains $k$ elements.

Doubling the table requires copying all $k$ existing elements, so the actual resizing cost is:

\[
k
\]

Observe that before expansion, the right half of the table contains:

\[
\frac{k}{2}
\]

elements, and by invariant each of them stores 2 credits. Thus the total saved credit is:

\[
2 \cdot \frac{k}{2} = k
\]

which is exactly enough to pay for copying all $k$ elements into the new table.

Therefore, every expansion is fully paid for using previously stored credits.

\subsubsection*{DELETE}

For each DELETE:

\begin{itemize}
    \item 1 unit pays for the immediate deletion,
    \item 1 unit is stored as credit in the newly empty slot.
\end{itemize}

Suppose a contraction occurs when the table has capacity $k$ and the number of elements becomes:

\[
\frac{k}{4}
\]

Then the table shrinks to capacity:

\[
\frac{k}{2}
\]

The contraction requires copying:

\[
\frac{k}{4}
\]

elements into the smaller table, so the resizing cost is:

\[
\frac{k}{4}
\]

Before shrinking, the left half of the table contains:

\[
\frac{k}{2}
\]

slots. Since only $\frac{k}{4}$ of the entire table are occupied, at least:

\[
\frac{k}{4}
\]

slots in the left half are empty.

Each empty slot stores 1 credit, so the total available credit is:

\[
\frac{k}{4}
\]

which is exactly enough to pay for moving the remaining elements during contraction.

Thus every contraction is also fully paid for by saved credits.

\subsubsection*{Total Cost}

Since every INSERT is charged 3 and every DELETE is charged 2, for any sequence of $n$ operations the total amortized cost is at most:

\[
3(\#\text{INSERT}) + 2(\#\text{DELETE})
\]

Since:

\[
\#\text{INSERT} + \#\text{DELETE} = n
\]

we obtain:

\[
3(\#\text{INSERT}) + 2(\#\text{DELETE}) \leq 3n
\]

Hence the total cost of $n$ operations is:

\[
O(n)
\]

and therefore the amortized cost per operation is:

\[
O(1)
\]


\subsection*{3. Potential Method}

Let:

\begin{itemize}
    \item $m$ = number of elements in the table,
    \item $T$ = table size (capacity).
\end{itemize}

Define the potential function:

\[
\Phi(T,m) =
\begin{cases}
2m - T & \text{if } m \geq \frac{T}{2} \\[6pt]
\frac{T}{2} - m & \text{if } m < \frac{T}{2}
\end{cases}
\]

The amortized cost is:

\[
\hat{c}_i = c_i + \Phi_i - \Phi_{i-1}
\]

For a normal INSERT without resizing:

\[
c_i = 1
\]

The change in potential is at most constant, so:

\[
\hat{c}_i = O(1)
\]

For an INSERT that triggers doubling:

\[
c_i = m + 1
\]

But the potential drops enough to pay for the copying cost, giving:

\[
\hat{c}_i = O(1)
\]

Similarly, for DELETE operations:

\begin{itemize}
    \item Normal DELETE has constant amortized cost.
    \item DELETE causing contraction also has constant amortized cost because the decrease in potential pays for copying.
\end{itemize}

Therefore:

\[
\hat{c}_i = O(1)
\]

for every operation, and for $n$ operations:

\[
\sum_{i=1}^{n} \hat{c}_i = O(n)
\]

\subsection*{Final Result}

Using aggregate, accounting, and potential methods:

\[
\boxed{\text{Total cost of } n \text{ INSERT and DELETE operations is } O(n)}
\]

Hence, the amortized cost per operation is:

\[
\boxed{O(1)}
\]

\section{Exercise 2}

\subsection*{a. SEARCH Operation}

There are at most

\[
k = \lceil \log(n+1) \rceil
\]

arrays, and array \(A_i\) has size \(2^i\).

To search for an element \(x\):

\begin{enumerate}
    \item For every nonempty array \(A_i\),
    perform binary search on \(A_i\).
    \item If \(x\) is found in any array, return TRUE.
    \item Otherwise, return FALSE.
\end{enumerate}

Since binary search on \(A_i\) takes

\[
O(\log 2^i) = O(i),
\]

the total worst-case running time is

\[
\sum_{i=0}^{k-1} O(i)
=
O\left(\sum_{i=0}^{k-1} i\right)
=
O(k^2).
\]

Since

\[
k = \Theta(\log n),
\]

we obtain

\[
O(k^2) = O((\log n)^2).
\]

Thus the worst case for SEARCH operation runs in

\[
\boxed{O((\log n)^2)}
\]

\subsection*{b. INSERT Operation}

To insert a new element \(x\):

\begin{enumerate}
    \item Create a new \textit{carry} array containing only \(x\).
    \item If \(A_0\) is empty, place the \textit{carry} there.
    \item Otherwise, merge the \textit{carry} with \(A_0\), producing a sorted array of size \(2\).
    \item If \(A_1\) is empty, place the merged array there.
    \item Otherwise, merge again and continue upward, exactly like carrying in binary addition.
\end{enumerate}

If arrays \(A_0, A_1, \dots, A_{t-1}\) are full and \(A_t\) is empty, then we perform merges of sizes

\[
1,2,4,\dots,2^{t-1}.
\]

The total cost is

\[
1 + 2 + 4 + \cdots + 2^{t-1}
=
2^t - 1.
\]

In the worst case, all arrays are full, so

\[
2^k - 1 = 2^{\log(n+1)} - 1 \approx 2^{\log n} = n^{\log2} = n
\]

and inserting one new element requires merging all elements.

Thus the worst-case running time is

\[
\boxed{O(n)}.
\]

\subsubsection*{Amortized Analysis: Aggregate Method}

Consider a sequence of \(n\) INSERT operations.

An array of size \(2^i\) is created only after every \(2^i\) insertions.
Whenever such an array is created, the merge cost is \(O(2^i)\).

Thus the total contribution of level \(i\) over \(n\) insertions is

\[
O\left(\frac{n}{2^i} \cdot 2^i\right)
=
O(n).
\]

There are at most $\lceil \log(n+1) \rceil$ levels, so the total cost of all insertions is

\[
O(n \log n).
\]

Therefore, the amortized cost per INSERT is

\[
\frac{O(n \log n)}{n}
=
O(\log n).
\]

Hence,

\[
\boxed{\text{Amortized INSERT cost } = O(\log n)}.
\]

\subsubsection*{Amortized Analysis: Accounting Method}

Charge each INSERT operation

\[
c = d \log n
\]

credits for a sufficiently large constant \(d\).

Whenever an element participates in a merge, it moves from an array of size \(2^i\) to an array of size \(2^{i+1}\).

Since the largest array has size at most \(n\), any element can move at most

\[
\log n
\]

times.

Each movement costs \(O(1)\) per element.

Thus each inserted element incurs a total future movement cost of at most

\[
O(\log n).
\]

Therefore charging \(O(\log n)\) credits per insertion is enough to pay for:

\begin{itemize}
    \item the initial insertion cost,
    \item all future merge costs involving that element.
\end{itemize}

Hence the amortized INSERT cost is

\[
\boxed{O(\log n)}.
\]

\end{document}
